% Section experiencia corta - DOCENCIA E INVESTIGACIÓN
\sectionTitle{Experiencias Profesionales}{\faSuitcase}

\begin{experiences}

\experience
{Ahora} {Docente de Metodología de la Investigación}{Corporación Educativa Universitaria}{Ayacucho}
{Enero 2026} {
	Acompañamiento integral en elaboración de monografías y proyectos de tesis con énfasis en rigor metodológico y ética investigativa.
	\begin{itemize}
		\item Asesoría personalizada en planteamiento de problemas de investigación y diseño metodológico
		\item Talleres de IA aplicada a la investigación: uso ético de ChatGPT, búsqueda académica y gestión bibliográfica
		\item Formación en normas APA 7, prevención de plagio académico y redacción científica
		\item Sesiones prácticas de análisis de datos con R/SPSS y elaboración de marco teórico
		\item Evaluación formativa mediante rúbricas y retroalimentación constructiva continua
	\end{itemize}
}
{\footnotesize{\emph{Herramientas:} Google Classroom, Turnitin, Zotero, Mendeley, R/SPSS, modalidad híbrida}}
\emptySeparator

\experience
{Diciembre 2025} {Docente de Metodología de la Investigación}{Corporación Educativa Universitaria}{Ayacucho}
{Enero 2025} {
	Formación en fundamentos metodológicos y diseño de proyectos de investigación científica.
	\begin{itemize}
		\item Enseñanza de epistemología, tipos de investigación y enfoques metodológicos
		\item Diseños cuantitativos, cualitativos y mixtos aplicados a ciencias sociales
		\item Elaboración, validación y confiabilización de instrumentos de recolección de datos
		\item Técnicas de muestreo probabilístico y no probabilístico
		\item Estrategias de análisis de datos: análisis estadístico y análisis de contenido
	\end{itemize}
}
{\footnotesize{\emph{Herramientas:} Moodle, Google Classroom, rúbricas analíticas, Mentimeter, Kahoot}}
\emptySeparator

\experience
{Julio 2024} {Docente de Gestión de Proyectos II}{Universidad Nacional Autónoma de Huanta}{Huanta}
{Abril 2024} {
	Docencia en curso de gestión de proyectos con enfoque cuantitativo y análisis de viabilidad.
	\begin{itemize}
		\item Diseño e implementación de proyectos aplicados con análisis costo-beneficio
		\item Evaluación económica y financiera de proyectos de inversión
		\item Elaboración de reportes técnicos y presentaciones ejecutivas
		\item Uso de herramientas de gestión: cronogramas, presupuestos, indicadores de desempeño
		\item Retroalimentación personalizada basada en competencias y rúbricas
	\end{itemize}
}
{\footnotesize{\emph{Tecnologías utilizadas:} Excel avanzado, MS Project, Trello, Google Workspace, LaTeX}}
\emptySeparator

%%%%%%%%%%%%%%%%%%%%%%%%%%%%%%%%%%%%%%%%
% EXPERIENCIAS ADICIONALES - Descomentar según relevancia
%%%%%%%%%%%%%%%%%%%%%%%%%%%%%%%%%%%%%%%%

%\experience
%{Agosto 2025} {Tutor de Tesis de Pregrado}{Universidad Nacional de San Cristóbal de Huamanga}{Ayacucho}
%{Enero 2024} {
%	Dirección y asesoría de tesis en economía con énfasis en investigación aplicada y rigor metodológico.
%	\begin{itemize}
%		\item Dirección de 5 tesis en economía del desarrollo, econometría aplicada y políticas públicas
%		\item Asesoría en diseño de investigación, construcción de marcos teóricos y revisión de literatura
%		\item Acompañamiento en análisis de datos con Stata/R: regresiones, series de tiempo, datos panel
%		\item Revisión exhaustiva de documentos siguiendo normas APA 7 y Vancouver
%		\item Preparación de tesistas para sustentación: elaboración de diapositivas y simulacros
%	\end{itemize}
%}
%{\footnotesize{\emph{Áreas temáticas:} Economía del desarrollo, econometría, evaluación de impacto, políticas públicas}}

%\experience
%{Diciembre 2023} {Coordinador de Investigación}{Instituto de Investigación}{Ayacucho}
%{Junio 2023} {
%	Coordinación de procesos investigativos institucionales y formación de semilleros de investigación.
%	\begin{itemize}
%		\item Gestión de convocatorias de proyectos de investigación docente y estudiantil
%		\item Organización de jornadas de investigación, seminarios metodológicos y talleres de publicación
%		\item Revisión técnica de proyectos y reportes finales de investigación
%		\item Seguimiento de cronogramas, presupuestos y entregables de proyectos de investigación
%		\item Coordinación con comités de ética y revisión por pares
%	\end{itemize}
%}
%{\footnotesize{\emph{Herramientas:} Sistema de gestión de proyectos institucional, OJS para revistas científicas}}

%\experience
%{Mayo 2023} {Facilitador de Talleres de Investigación}{Diversos}{Ayacucho}
%{Enero 2022} {
%	Facilitación de talleres especializados en metodología de investigación para profesionales y estudiantes.
%	\begin{itemize}
%		\item Taller "Diseño de Investigación Cuantitativa" - 40 horas
%		\item Taller "Análisis de Datos con R para Investigadores" - 32 horas
%		\item Taller "Redacción Científica y Publicación en Revistas Indexadas" - 24 horas
%		\item Taller "Elaboración de Instrumentos y Validación de Contenido" - 16 horas
%		\item Metodología participativa con ejercicios prácticos y casos aplicados
%	\end{itemize}
%}
%{\footnotesize{\emph{Participantes:} Docentes universitarios, investigadores, profesionales de la salud y educación}}

%\experience
%{Diciembre 2022} {Investigador Junior}{Centro de Investigación Económica y Social}{Ayacucho}
%{Marzo 2022} {
%	Participación en proyectos de investigación económica y social aplicada.
%	\begin{itemize}
%		\item Revisión sistemática de literatura en bases como Scopus, Web of Science, EBSCO
%		\item Procesamiento y análisis de encuestas y bases de datos secundarias (ENAHO, ENE)
%		\item Análisis estadístico con Stata: estadística descriptiva, correlaciones, regresiones
%		\item Elaboración de secciones metodológicas y de resultados para reportes técnicos
%		\item Apoyo en redacción de artículos científicos y working papers
%	\end{itemize}
%}
%{\footnotesize{\emph{Proyectos:} Evaluación de programas sociales, análisis de mercado laboral regional}}

%\experience
%{Agosto 2021} {Asistente de Cátedra en Estadística}{Universidad Nacional de San Cristóbal de Huamanga}{Ayacucho}
%{Marzo 2021} {
%	Apoyo docente en curso de Estadística para economistas.
%	\begin{itemize}
%		\item Dictado de clases prácticas y resolución de ejercicios
%		\item Asesoría personalizada a estudiantes en laboratorio de cómputo
%		\item Calificación de evaluaciones y retroalimentación formativa
%		\item Desarrollo de material didáctico: ejercicios, casos y tutoriales
%		\item Apoyo en elaboración de exámenes y rúbricas de evaluación
%	\end{itemize}
%}
%{\footnotesize{\emph{Software utilizado:} Excel, Stata, R (nivel introductorio para estudiantes)}}

\end{experiences}
